\uuid{87v9}
\titre{Circuit RLC en régime critique}

\niveau{M1}
\module{Analyse}
\chapitre{Systèmes linéaires causals}
\sousChapitre{Réponse impulsionnelle d'un circuit RLC}
\theme{Convolution et signal de sortie}
\auteur{Quentin Liard}
\datecreate{ 2026-09-03}
\organisation{CentraleSupélec}
\difficulte{4}

\contenu{

\texte{
On considère un circuit RLC dont la tension de sortie $v$ est liée à la tension d'entrée $u$ par
$$
\frac{L}{R}v''(t)+v'(t)+\frac{v(t)}{RC}=u'(t).
$$
On suppose $R,L,C>0$ et
$$
R=2\sqrt{\frac{L}{C}}.
$$
}

\begin{enumerate}

\item \question{Déterminer toutes les solutions de l'équation différentielle
$$
\frac{L}{R}z''(t)+z'(t)+\frac{z(t)}{RC}=0.
$$}

\reponse{
L'hypothèse $R=2\sqrt{L/C}$ donne
$$
\frac{L}{R}=\frac{RC}{4}.
$$
L'équation devient donc
$$
\frac{RC}{4}z''+z'+\frac{1}{RC}z=0,
$$
soit, après multiplication par $4/(RC)$,
$$
z''+\frac{4}{RC}z'+\frac{4}{R^2C^2}z=0.
$$
L'équation caractéristique est
$$
r^2+\frac{4}{RC}r+\frac{4}{R^2C^2}=0
\quad\Longleftrightarrow\quad
\left(r+\frac{2}{RC}\right)^2=0.
$$
Il existe une racine double $-2/(RC)$. Les solutions sont donc
$$
\boxed{z(t)=(\alpha+\beta t)e^{-\frac{2t}{RC}},\qquad \alpha,\beta\in\mathbb{R}}.
$$
}

\item \question{En choisissant des conditions initiales convenables, déterminer une solution $G\in\mathcal{D}'_+$ de
$$
\frac{L}{R}G''+G'+\frac{G}{RC}=\delta.
$$}

\reponse{
Choisissons la solution homogène $z$ vérifiant
$$
z(0)=0,
\qquad
z'(0)=1.
$$
D'après la question précédente,
$$
z(t)=t e^{-\frac{2t}{RC}}.
$$
Pour une telle fonction,
$$
(Hz)''=Hz''+\delta.
$$
Si l'on note
$$
D=\frac{L}{R}\frac{d^2}{dt^2}+\frac{d}{dt}+\frac{1}{RC},
$$
alors
$$
D(Hz)=\frac{L}{R}\delta.
$$
Il suffit donc de poser
$$
G=\frac{R}{L}Hz.
$$
Or, sous l'hypothèse imposée,
$$
\frac{R}{L}=\frac{4}{RC}.
$$
Ainsi,
$$
\boxed{G(t)=\frac{4}{RC}H(t)t e^{-\frac{2t}{RC}}}.
$$
}

\item \question{En déduire une relation de la forme
$$
v=g\ast u,
$$
et préciser la fonction $g$.}

\reponse{
Sous l'hypothèse que les convolutions considérées sont définies, notamment pour des signaux causaux, une solution est
$$
v=G\ast u'.
$$
Comme la dérivation commute à la convolution,
$$
v=G'\ast u.
$$
Écrivons
$$
G(t)=\frac{4}{RC}H(t)q(t),
\qquad
q(t)=t e^{-\frac{2t}{RC}}.
$$
Comme $q(0)=0$,
$$
G'(t)=\frac{4}{RC}H(t)q'(t).
$$
Or
$$
q'(t)=\left(1-\frac{2t}{RC}\right)e^{-\frac{2t}{RC}}.
$$
On obtient donc
$$
\boxed{v=g\ast u},
$$
avec
$$
\boxed{
g(t)=\frac{4}{RC}
\left(1-\frac{2t}{RC}\right)
H(t)e^{-\frac{2t}{RC}}
}.
$$
}

\end{enumerate}

}
